\chapter*{Preface}
\addcontentsline{toc}{chapter}{Preface}
\markboth{Preface}{Preface}

Every time you unlock a phone, stream a video, pay with a card or glance at a map, you
rely on a chain of ideas that runs back through the space program, the transistor, radar,
the vacuum tube, the telephone and the telegraph. Each link in that chain was forged by
engineers solving a concrete problem---how to send more messages over the same wire, how
to hear a ship beyond the horizon, how to fit a thousand conversations into a slice of
spectrum---and each solution left behind a principle that is still at work inside
today's modems. This book is an attempt to hand those principles to you the way a senior
engineer hands them to an apprentice: with the mathematics that makes them exact, the
history that explains why they look the way they do, the hardware that makes them real,
and a large number of simulations that let you break them and watch what happens.

A communication engineer's knowledge is peculiar. It is deeply mathematical---the
matched filter, the Shannon limit and the Viterbi algorithm are theorems---and yet it is
also a craft, full of rules of thumb about loop bandwidths, back-off, guard bands and
link margins that you only learn by building things that fail. I have tried to give both
halves equal weight. When a result is exact, it is derived. When it is a rule of thumb,
the text says so and explains where it comes from, so that you will know when it stops
applying.

\section*{Who this book is for}
The book is written for engineers and advanced students in electrical and computer
engineering. It assumes you are comfortable with calculus, linear algebra, basic
probability and a programming language; it does not assume a prior course in
communications or signal processing. Early chapters build the signal-processing and
probability foundations quickly but completely, so that a computer engineer can follow
the whole arc, while the boxed discussions and the problem sets push well into material
normally reserved for graduate courses.

\section*{How the book is organised}
\begin{description}[leftmargin=0pt]
\item[Part I, Foundations,] opens with the history of telecommunication and then develops
signals and spectra, random processes and noise, the classic analog systems---AM, FM,
single sideband, the superheterodyne receiver, analog television---and the step into
digital telephony through sampling, quantization and PCM.
\item[Part II, Digital Signal Processing for Radio,] covers the filters, multirate
structures, oscillators and converters that every modern radio is built from, and
dissects the transceiver and the software-defined radio.
\item[Part III, Digital Transmission,] is the heart of physical-layer engineering:
pulse shaping, modulation and detection, synchronization, propagation and equalization.
\item[Part IV, Information and Coding,] gives Shannon's theory and the codes that
approach its limits, from Hamming and Reed--Solomon to turbo, LDPC and polar codes,
together with the source coding that compresses voice, audio and images.
\item[Part V, Systems,] assembles everything into real systems: OFDM, spread spectrum and
GPS, multiple antennas, multiple access and the cellular concept, every cellular
generation from AMPS to 5G, Wi-Fi and the Internet of Things, satellite links from
Telstar to Starlink, wireline and fibre-optic systems, and the research frontier of 6G.
\end{description}

A one-semester introductory course can cover Chapters 1--5 and 8--11. A second course
can take Chapters 12--19. Chapters 20--25 read well on their own as a systems course or
as professional reference.

\section*{Features}
Six kinds of boxes appear throughout the text.
\begin{history}[]The people and events behind an idea, and why the solution took the form it did.\end{history}
\begin{inpractice}[]Where the idea lives in deployed equipment and standards, with real numbers.\end{inpractice}
\begin{keyidea}[]A principle worth remembering long after the derivation is forgotten.\end{keyidea}
\begin{pitfall}[]A mistake that experienced engineers still make.\end{pitfall}
\begin{worked}[]The arithmetic of a real design, carried through to a number you can check.\end{worked}
\begin{labbox}[]A Python simulation in the companion labs that lets you explore the topic.\end{labbox}
Every data figure in the book was produced by Python code in the companion repository,
usually by the same library functions the labs use, so any plot can be regenerated,
modified and extended. Worked examples show the arithmetic of real designs, and every
chapter ends with problems graded from routine to open-ended and with annotated further
reading.

\section*{The companion labs}
The \texttt{moderncomms-labs} repository contains the \texttt{commlib} signal-processing
library, thirty-six interactive Jupyter labs (one or more for every chapter), the scripts that generate every
figure in this book, and GNU Radio flowgraphs for the Ettus USRP B200 software-defined
radio. The labs run entirely in simulation; the hardware examples are optional but
strongly encouraged. Appendix~B explains how to install and use them.

\section*{A note on numbers and standards}
Standards evolve. Figures for 5G, Wi-Fi~7, satellite constellations and the 6G study
reflect publicly available information as of September~2026. Where a number is an
engineering approximation rather than a specification value, the text says so.
